\section{Results and Discussion}
\label{sec:results}

\subsection{Evidence from the hourly source study}

Under the existing 72-hour hourly protocol, TimesNet achieved a macro-RMSE of
101.04~W/m$^2$, compared with 104.68~W/m$^2$ for the LSTM, a relative
reduction of 3.48\%. TimesNet also obtained lower RMSE than the LSTM at the
24- and 48-hour prefixes and lower RMSE than both temporal baselines at all
three prefixes. These values describe the hourly source experiment only.

\subsection{Evidence from the monthly source study}

In the existing one-month-ahead experiment, the consolidated DilatedRNN
forecast achieved a macro-MAE of 13.42~W/m$^2$, an RMSE of 16.20~W/m$^2$,
and an $R^2$ of 0.898. It ranked first among five evaluated methods by
macro-MAE, but led at only four of the ten locations. Local MAE ranged from
7.45 to 20.92~W/m$^2$. These values describe the monthly source experiment
only and are not numerically compared with the hourly TimesNet errors.

\subsection{Matched multi-resolution comparison}

\DraftNote{Insert the daily and matched TimesNet--DilatedRNN tables generated
from the new experiment artifacts. Report one overall table by resolution and
horizon, followed by location-level comparisons.}

The final analysis will distinguish three questions: which model minimizes the
declared primary metric on average, how often that model wins at matched test
origins, and whether its advantage is consistent across locations. A method
will not be described as universally superior when the location-level evidence
shows losses.

\subsection{Favorable and unfavorable cases}

\DraftNote{Create one table for the strongest relative gains and another for
the largest deficits of the proposed method. Include location, date or origin,
horizon, reference GHI, both forecasts, absolute errors, and confirmed weather
context when available.}

\subsection{Geographic and climatic interpretation}

Climate zones and factory locations will be used as descriptive context. With
only ten locations, apparent differences between zones cannot by themselves
establish a causal climatic effect. Sites in the same country may remain far
apart and occupy distinct regimes, while nearby sites may share correlated
weather. The discussion will therefore combine the map, zone classification,
and error profiles without overstating statistical evidence.

